\subsection{Computation of composition products by shifting resolutions}

Let

\[
0 \rightarrow \mathrm{M} \xrightarrow {f} \mathrm{K} _ {n} \rightarrow \mathrm{K} _ {n - 1} \rightarrow \dots \rightarrow \mathrm{K} _ {1} \xrightarrow {g} \mathrm{M} ^ {\prime} \rightarrow 0\tag{16}
\]

be an exact sequence of left A-modules and \(\theta\in\mathrm{Ext}_{\mathbf{A}}^{n}(\mathbf{M}^{\prime},\mathbf{M})\) the associated class.

Let \(a:(\mathbf{R},d)\to \mathbf{M}\) be a left resolution of \(\mathbf{M}\) ; we thus have an exact sequence

\[
\rightarrow \mathrm{R} _ {k} \xrightarrow {d _ {k}} \mathrm{R} _ {k - 1} \rightarrow \dots \xrightarrow {d _ {1}} \mathrm{R} _ {0} \xrightarrow {a _ {0}} \mathrm{M} \rightarrow 0.
\]

and by translation of n (X, p.~26) an exact sequence

\[
\rightarrow \mathrm{R} _ {k} \xrightarrow {(- 1) ^ {n} d _ {k}} \mathrm{R} _ {k - 1} \rightarrow \dots \xrightarrow {(- 1) ^ {n} d _ {1}} \mathrm{R} _ {0} \xrightarrow {(- 1) ^ {n} a _ {0}} \mathrm{M} \rightarrow 0.\tag{17}
\]

From (16) and (17) we deduce an exact sequence

\[
\rightarrow \mathrm{R} _ {k} \xrightarrow {(- 1) ^ {n} d _ {k}} \mathrm{R} _ {k - 1} \rightarrow \dots \xrightarrow {(- 1) ^ {n} d _ {1}} \mathrm{R} _ {0} \xrightarrow {(- 1) ^ {n} f \circ a _ {0}} \mathrm{K} _ {n} \rightarrow \mathrm{K} _ {n - 1} \rightarrow \dots \rightarrow \mathrm{K} _ {1} \rightarrow \mathrm{M} ^ {\prime} \rightarrow 0
\]

whence a resolution \(\mathbf{R}'\) of \(\mathbf{M}'\) ; denote by \(\varphi: \mathbf{R}' \to \mathbf{R}(-n)\) the morphism such that \(\varphi_k = 1_{\mathbf{R}_{k-n}}\) for \(k \geqslant n\) .

If N is a left A-module and P is a right A-module, we therefore have homomorphisms

\[
\mathrm{H} \left(1 _ {\mathrm{P}} \otimes \varphi\right): \mathrm{H} \left(\mathrm{P} \otimes_ {\mathrm{A}} \mathrm{R} ^ {\prime}\right)\rightarrow \mathrm{H} \left(\mathrm{P} \otimes_ {\mathrm{A}} \mathrm{R}\right) (- n)
\]

\[
\mathrm{H} \left(\operatorname{Homgr} _ {\mathbf {A}} (\varphi , 1 _ {\mathbf {N}})\right): \mathrm{H} \left(\operatorname{Homgr} _ {\mathbf {A}} (\mathrm{R}, \mathrm{N})\right) (n) \rightarrow \mathrm{H} \left(\operatorname{Homgr} _ {\mathbf {A}} \left(\mathrm{R} ^ {\prime}, \mathrm{N}\right)\right).
\]

Let k be an integer.

PROPOSITION 8. --- a) The following diagram, where \(h_{\theta}(\alpha) = \theta \circ \alpha\) , is commutative

\[
\begin{array}{c} \operatorname{Tor} _ {k + n} ^ {\mathrm{A}} (\mathrm{P}, \mathrm{M} ^ {\prime}) \xrightarrow {h _ {\theta}} \operatorname{Tor} _ {k} ^ {\mathrm{A}} (\mathrm{P}, \mathrm{M}) \\ \psi_ {k + n} (\mathrm{P}, \mathrm{R} ^ {\prime}) \Big | _ {\downarrow} \quad \Big | _ {\downarrow} \psi_ {k} (\mathrm{P}, \mathrm{R}) \\ \mathrm{H} _ {k + n} (\mathrm{P} \otimes_ {\mathrm{A}} \mathrm{R} ^ {\prime}) \xrightarrow {\mathrm{H} _ {k + n} (1 \otimes \varphi)} \mathrm{H} _ {k} (\mathrm{P} \otimes_ {\mathrm{A}} \mathrm{R}) \end{array}
\]

\begin{enumerate}
\def\labelenumi{\alph{enumi})}
\setcounter{enumi}{1}
\tightlist
\item
  The following diagram, where \(\delta_{\theta}(\beta) = \beta \circ \theta\) , is commutative
\end{enumerate}

\[
\begin{array}{c} \mathrm{H} ^ {k} (\text {Homgr} _ {\mathrm{A}} (\mathrm{R}, \mathrm{N})) \xrightarrow {\mathrm{H} ^ {k + n} (\text {Homgr} _ {\mathrm{A}} (\varphi , 1))} \mathrm{H} ^ {k + n} (\text {Homgr} _ {\mathrm{A}} (\mathrm{R} ^ {\prime}, \mathrm{N})) \\ \varphi^ {k} (\mathrm{R}, \mathrm{N}) \Biggl \downarrow \\ \operatorname{Ext} _ {\mathrm{A}} ^ {k} (\mathrm{M}, \mathrm{N}) \xrightarrow {\delta_ {\theta}} \operatorname{Ext} _ {\mathrm{A}} ^ {k + n} (\mathrm{M} ^ {\prime}, \mathrm{N}). \end{array}
\]

Let \(\alpha : L(M) \to R\) be a morphism of complexes such that \(a \circ \alpha = p_{M}\) and let

\[
\begin{array}{c} \mathrm{L} _ {n} (\mathrm{M} ^ {\prime}) \longrightarrow \mathrm{L} _ {n - 1} (\mathrm{M} ^ {\prime}) \longrightarrow \dots \longrightarrow \mathrm{L} _ {0} (\mathrm{M} ^ {\prime}) \longrightarrow \mathrm{M} ^ {\prime} \longrightarrow 0 \\ u _ {n} \Big \downarrow \qquad \qquad \qquad \Big \downarrow u _ {n - 1} \qquad \qquad \qquad \qquad \qquad \qquad \qquad \qquad \qquad \qquad \qquad \qquad \qquad \qquad \qquad \qquad \qquad \qquad \qquad \qquad \qquad \qquad \qquad \qquad \qquad \qquad \qquad \qquad \qquad \qquad \qquad \qquad \qquad \qquad \\ 0 \longrightarrow \mathrm{M} \xrightarrow {f} \mathrm{K} _ {n} \longrightarrow \dots \longrightarrow \mathrm{K} _ {1} \longrightarrow \mathrm{M} ^ {\prime} \longrightarrow 0 \end{array}
\]

a commutative diagram; let us choose a homomorphism \(v_{n}: \mathrm{L}_{n}(\mathrm{M}^{\prime}) \to \mathrm{L}_{0}(\mathrm{M})\) such that \(p_{\mathbf{M}} \circ v_{n} = (-1)^{n} u_{n}\) ; by X, p.~47, prop. 1, a), \(v_{n}\) extends to a morphism of complexes \(v: \mathrm{L}(\mathrm{M}^{\prime}) \to \mathrm{L}(\mathrm{M})(-n)\) , and \(\theta\) is the image under the canonical isomorphism \(\mathrm{H}^{n}(\mathrm{Homgr}_{\mathbf{A}}(\mathrm{L}(\mathrm{M}^{\prime}), \mathrm{L}(\mathrm{M}))) \to \mathrm{Ext}_{\mathbf{A}}^{n}(\mathbf{M}^{\prime}, \mathbf{M})\) of the class of \(v\) (X, p.~117, remark 1). One defines a morphism of complexes \(\beta: \mathrm{L}(\mathrm{M}^{\prime}) \to \mathrm{R}^{\prime}\) by \(\beta_{p} = u_{p}\) for \(p \leqslant n - 1\) , \(\beta_{p} = \alpha_{p-n} \circ v_{p}\) for \(p \geqslant n\) , and we have

\[
\varphi \circ \beta = \alpha (- n) \circ v.
\]

On the other hand, by definition of the maps \(\varphi\) and \(\psi\) , we have

\[
\begin{array}{r l r} & {\psi_ {k} (\mathrm{P}, \mathrm{R}) = \mathrm{H} _ {k} (p _ {\mathrm{P}} \otimes \alpha),} & {\varphi^ {k} (\mathrm{R}, \mathrm{N}) = \mathrm{H} ^ {k} (\mathrm{Homgr} _ {\mathrm{A}} (\alpha , e _ {\mathrm{N}})),} \\ & {\psi_ {k + n} (\mathrm{P}, \mathrm{R} ^ {\prime}) = \mathrm{H} _ {k + n} (p _ {\mathrm{P}} \otimes \beta),} & {\varphi^ {k + n} (\mathrm{R} ^ {\prime}, \tilde {\mathrm{N}}) = \mathrm{H} ^ {k + n} (\mathrm{Homgr} _ {\mathrm{A}} (\beta , e _ {\mathrm{N}})).} \end{array}
\]

Finally, by definition of the composition product, one has

\[
h _ {\theta} = \mathrm{H} (1 _ {L (\mathbb {P})} \otimes v), \quad \delta_ {\theta} = \mathrm{H} (\operatorname{Homgr} _ {\mathbf {A}} (v, 1 _ {I (\mathbb {N})}))
\]

Consequently, one has the equalities

\[
\begin{array}{r l} \psi_ {k} (\mathrm{P}, \mathrm{R}) \circ h _ {\theta} & = \mathrm{H} _ {k} (p _ {\mathrm{P}} \otimes \alpha) \circ \mathrm{H} _ {k} (1 _ {\mathrm{L(P)}} \otimes v) = \mathrm{H} _ {k} (p _ {\mathrm{P}} \otimes (\alpha \circ v)) = \mathrm{H} _ {k + n} (p _ {\mathrm{P}} \otimes (\varphi \circ \beta)) \\ & = \mathrm{H} _ {k + n} (1 \otimes \varphi) \circ \mathrm{H} _ {k + n} (p _ {\mathrm{P}} \circ \beta) = \mathrm{H} _ {k + n} (1 \otimes \varphi) \circ \psi_ {k + n} (\mathrm{P}, \mathrm{R} ^ {\prime}), \end{array}
\]

whence a); the proof of b) is analogous.

Remark. --- By the commutation isomorphisms, we deduce from \(a\) ) an analogous statement in the case of an exact sequence (16) of right A-modules.

Now let \(b: \mathbf{M}' \to \mathbf{E}'\) be a right resolution of \(\mathbf{M}'\) ; we thus have an exact sequence

\[
0 \rightarrow \mathbf {M} ^ {\prime} \xrightarrow {\dot {b} ^ {0}} \mathrm{E} ^ {\prime 0} \xrightarrow {\delta^ {0}} \mathrm{E} ^ {\prime 1} \rightarrow \dots \rightarrow \mathrm{E} ^ {\prime k} \xrightarrow {\delta^ {k}} \mathrm{E} ^ {\prime k + 1}
\]

whence an exact sequence

\[
0 \rightarrow \mathrm{M} \xrightarrow {f} \mathrm{K} _ {n} \rightarrow \mathrm{K} _ {n - 1} \rightarrow \dots \rightarrow \mathrm{K} _ {1} \xrightarrow {(- 1) ^ {n} b ^ {0} \circ g} \mathrm{E} ^ {\prime 0} \xrightarrow {(- 1) ^ {n} \delta^ {0}} \mathrm{E} ^ {\prime 1} \rightarrow \dots
\]

corresponding to a right resolution E of M; denote \(\sigma : E'(n) \to E\) the morphism such that \(\sigma^k = 1_{E'k - n}\) for \(k \geqslant n\) We therefore have homomorphisms

\[
\mathrm{H} \left(\operatorname{Homgr} _ {\mathrm{A}} \left(1 _ {\mathrm{N}}, \sigma\right)\right): \mathrm{H} \left(\operatorname{Homgr} _ {\mathrm{A}} \left(\mathrm{N}, \mathrm{E} ^ {\prime}\right)\right) (n) \rightarrow \mathrm{H} \left(\operatorname{Homgr} _ {\mathrm{A}} (\mathrm{N}, \mathrm{E})\right).
\]

PROPOSITION 9. --- The following diagram, where \(\gamma_{\theta}(\alpha) = \theta \circ \alpha\) , is commutative:

\[
\begin{array}{c} \mathrm{H} ^ {k} (\text {Homgr} _ {\mathrm{A}} (\mathrm{N}, \mathrm{E} ^ {\prime})) \xrightarrow {\mathrm{H} ^ {k + n} (\text {Homgr} _ {\mathrm{A}} (1 _ {\mathrm{N}} , \sigma))} \mathrm{H} ^ {k + n} (\text {Homgr} _ {\mathrm{A}} (\mathrm{N}, \mathrm{E})) \\ \varphi^ {k} (\mathrm{N}, \mathrm{E} ^ {\prime}) \Big \downarrow \\ \operatorname{Ext} _ {\mathrm{A}} ^ {k} (\mathrm{N}, \mathrm{M} ^ {\prime}) \xrightarrow {\gamma_ {\theta}} \operatorname{Ext} _ {\mathrm{A}} ^ {k + n} (\mathrm{N}, \mathrm{M}). \end{array}
\]

This is proved in a manner analogous to prop. 8.

